\documentclass{article}
\usepackage{amsmath,amssymb}
\title{Shared-state chronological assembly and native synchronization}
\date{15 November 2026}
\begin{document}\maketitle
\paragraph{Status.} General construction independently reviewed under the stated hypotheses. It addresses an actual sharing law missing from independent-state tensor assembly. The input is an ordered chain with one retained state throughout. Its result is not an unqualified parallel sharing or feedback-fusion law.

\paragraph{Independent categorical input.} Fix the original state bicategory $\mathcal S$, process bicategory $\mathcal K$, and pure-action pseudofunctors $E_A:\mathcal S\to\mathcal K$ for external interfaces $A$. Fix a finite interface list $\mathbf A=(A_0,\ldots,A_m)$ with $m\geq1$. Define a common-state chain bicategory $C_{\mathbf A}$ explicitly. Its objects consist of one state $X$ and bodies
\[
 f_i:E_{A_{i-1}}X\longrightarrow E_{A_i}X\qquad(1\leq i\leq m).
\]
An arrow is one state map $h:X\to Y$ and, for every $i$, an actual native square
\[
 \alpha_i:E_{A_i}h\, f_i\Rightarrow g_i E_{A_{i-1}}h.
\tag{1}
\]
A two-cell is one state cell $\sigma:h\Rightarrow k$ satisfying the native square equation for every $i$. Composition and identity paste the corresponding squares and use the same state constraints. Bicategory equations hold componentwise after native admission. This explicit presentation avoids silently replacing exact common state and common witnesses by independent states, or by an unspecified bicategorical pullback.

\paragraph{Chronological evaluation.} Choose parentheses for the original process composite and define
\[
 C(X,\mathbf f)=\bigl(X,f_m\cdots f_1\bigr)
 \in D(A_0,A_m).
\tag{2}
\]
Keep $h$ on arrows and paste (1) from the last step toward the first:
\[
 E_{A_m}h\,f_m\cdots f_1
 \Rightarrow g_m E_{A_{m-1}}h\,f_{m-1}\cdots f_1
 \Rightarrow\cdots\Rightarrow
 g_m\cdots g_1 E_{A_0}h.
\tag{3}
\]
This is the chronological original-body composite; the reverse order of moving the boundary in (3) does not reverse execution. Retain $\sigma$ on two-cells. Substituting its individual native equations into (3) proves its output equation. Interchange proves preservation of arrow composition, with the original pseudofunctor frames supplying the comparison. Thus $C$ is a pseudofunctor over $1_{\mathcal S}$. If every $\alpha_i$ is invertible, so is (3). No pure tensor, diagonal, state copying, or inverse effectful interchanger occurs in this construction.

Associative regrouping of the finite chain gives native equivalences over the identity state map, with body the actual process associator. Their naturality is process associator naturality on the squares (1). Their coherence equations are the original process pentagon and unit equations. This gives coherent concatenation of chains when they have the same state and interface at the join. It does not equip independently varying-state arrows with a new total binary composition: the shared state and its common witness are genuine domain constraints.

\paragraph{Created synchronization cones.} Suppose the selected incoming lax bilimits of the underlying state diagrams exist in $\mathcal S$, and each output pseudofunctor $E_{A_i}$ for $1\leq i\leq m$ preserves them. For a pseudodiagram $F:J\to C_{\mathbf A}$, choose its underlying state cone $L\to X_j$. The native dialgebra cone-creation proof applies to each component body using this same state cone. The required output preservation is available separately for every component. It constructs one body $\bar f_i:E_{A_{i-1}}L\to E_{A_i}L$ and pseudo native projection square for each $i$. Together these form a chain cone on state $L$.

The universal property is likewise simultaneous: a source chain cone first yields its unique-up-to-equivalence state cone map; the individual native lift equations determine the component process data with that same map. A two-cell is admitted if and only if all component equations hold. Hence the multi-operation cone-creation argument gives the selected lax bilimit in $C_{\mathbf A}$. This is the previously established multi-symbol form of native cone creation, applied to a common-state chain, not a new assumption that independent cone choices happen to agree.

The target $D(A_0,A_m)$ has its native cones created over the same state cones, since $E_{A_m}$ preserves them. Chronological evaluation (2) is over the identity state base and preserves pseudo arrows by (3). The existing native preservation criterion therefore gives a cone-compatible equivalence
\[
 C\,L_J^{C_{\mathbf A}}F\ \simeq\
 L_J^{D(A_0,A_m)}(CF).
\tag{4}
\]
This is strong synchronization compatibility for same-state chronological assembly. It does not require a tensor of states to preserve a shared-index limit. Naturality includes full native diagram arrows and modifications. For lax-body diagram maps it follows from the actual cone transport of the component squares, with no inverse lax naturality cell. Shape restriction and its directed two-cells are compatible by the same cone projections.

Regrouping the chain commutes coherently with (4). The two composite comparison arrows, with the original process associator inserted, have projected cones related by its native naturality on the same common-state cone legs. Fullness supplies their invertible native two-cell. Equations between the resulting two-cell pastes follow by faithfulness from the original process coherence equations; no further comparison between modifications is introduced. Endpoint changes are governed by the previously derived relative pure-endpoint comparisons, whenever the original action package supplies their typed tightening square.

\paragraph{Feedback meaning and exact distinction.} For an exposed body on a retained state $X$, closing once gives the machine $(X,f)$. Closing the composite chain once therefore gives $(X,f_m\cdots f_1)$. Its later body sees the state produced by the earlier body. Independently closing each body and composing those machines instead retains separate state coordinates, normally $X^{\otimes m}$, with the existing held-body chronological semantics. These are different operations. Neither (2) nor (4) identifies them. A state reparameterization could relate them only if its actual native square is supplied and verified.

In particular, no concurrent shared-state operation is inferred from (2). Two effects on the same state may fail to commute or may require a specified joint transition law. An independent-state parallel operation does not select that law. The present theorem supplies ordered shared-state interference with exact side-effect forwarding, and its compatibility with native synchronization; it leaves scheduling and simultaneous interaction to actual admitted process data.

\paragraph{Sanity check, not a premise.} In the locally discrete Cartesian process category of finite sets take a bit state $S$. Let $f:S\to S$ flip the bit and let $g:S\to S\times S$ return its state as output while retaining it. Shared-state chronological closure outputs the flipped bit. Independent closures with initial pair $(s,s)$ output the original second bit. The diagonal $S\to S\times S$ therefore fails the required transition square between these machines: it would equate different outputs. This witnesses concretely why one must not silently replace early independent closure by late shared closure. The categorical theorem above does not rely on this example or assume Cartesian state.

\paragraph{Scope and attribution.} Chronological pasting and common-state native dialgebras are established categorical mechanisms. The claim is their precise integration with the original retained-state presentation and the earlier cone-creation criterion. No novelty is claimed for composition or the criterion. This is a finite-chain theorem under explicit output-preservation hypotheses, not a completion of arbitrary instantaneous feedback, arbitrary scheduler semantics, or an assertion that every desired joint physical law exists.
\end{document}
